%============================================================================
% LICENSING NOTICE
%============================================================================
% This WIP file, upon integration into guide-v.tex, becomes part of the
% TMS/DMS/CMS Usage Guide and is released under CC BY-NC 4.0.
% For details, see LICENSE in the repository root.
%============================================================================

%============================================================================
% FALCON BMS TMS/DMS/CMS HOTAS GUIDE
% WIP FILE TEMPLATE V1.0 – FINAL (Briefing v0.2.0.1 + TOC Fix)
%============================================================================

% IMPORTANTE: Este é um TEMPLATE padronizado para arquivos WIP (Work-In-Progress)
% que serão integrados ao guide.tex.

% Nomenclatura: Siga wip-naming-v1.4
% Padrão: chapter-C{N}-{titulo}-{status}-{data}.tex
% Exemplo: chapter-C2-hotas-fundamentals-dev-2026-02-06.tex

% Status: dev | review | final | approved | deprecated
% Localização: WIP/ (ativo) | ARCHIVE/ (aprovado/descartado)

%============================================================================
% PREAMBLE COMPLETO – TMS/DMS/CMS Usage Guide for Falcon BMS 4.38.1
% Gerado: 29 January 2026
% Status: Novo preâmbulo (report + twoside + titlesec + fancyhdr melhorado +
%soul + caption + new hotastable environment)
%============================================================================

\documentclass[11pt, a4paper, twoside]{report}

%--------------------------------------------------------------------------
% BASIC ENCODING AND LANGUAGE
%--------------------------------------------------------------------------

\usepackage[utf8]{inputenc}
\usepackage[T1]{fontenc}
\usepackage[english]{babel}

%--------------------------------------------------------------------------
% FONTS AND MICROTYPOGRAPHY
%--------------------------------------------------------------------------

\usepackage{lmodern}
\usepackage{microtype}

%--------------------------------------------------------------------------
% PAGE GEOMETRY AND LAYOUT
%--------------------------------------------------------------------------

\usepackage{geometry}
\geometry{a4paper, left=2.0cm, right=2.0cm, top=2.5cm, bottom=2.5cm}
\usepackage{setspace}
\onehalfspacing

% --------------------------------------------------------------------------
% COLORS AND LINKS
% --------------------------------------------------------------------------

\usepackage[table]{xcolor}
\definecolor{linkblue}{HTML}{004488}
\definecolor{linkred}{HTML}{882222}
\definecolor{headerblue}{HTML}{003366}
\definecolor{rowgray}{HTML}{F5F5F5}
\definecolor{subheadgray}{HTML}{E0E0E0}

\usepackage{soul}
\usepackage[pdfencoding=auto, psdextra, colorlinks=true, linkcolor=linkblue, citecolor=linkred, urlcolor=linkblue, breaklinks=true]{hyperref}
\usepackage{bookmark}

% ============================================================================
% CAPTION SETUP BY TYPE
% ============================================================================

\usepackage{caption}

% GLOBAL PATTERN
\captionsetup{font=small, labelfont=bf, justification=centering, singlelinecheck=true}
% FIGURES
\captionsetup[figure]{font=footnotesize, labelfont=bf, justification=centering, singlelinecheck=true}
% % TABLE/LONGTABLE/hotastable:
\captionsetup[table]{font=small, labelfont=bf, justification=centering, singlelinecheck=true}

%--------------------------------------------------------------------------
% HEADERS AND FOOTERS (IMPROVED for report + twoside)
%--------------------------------------------------------------------------

\usepackage{fancyhdr}
\setlength{\headheight}{25pt}                    % Increased (was 15pt) for long names
\pagestyle{fancy}
\fancyhf{}                                        % Clear all
\fancyhead[LO,RE]{\small\textit{\leftmark}}     % Outer edge (odd left, even right): chapter name
\fancyhead[RO,LE]{\small\thepage}               % Inner edge (odd right, even left): page number
\fancyfoot{}                                      % No footer (page number in header)
\renewcommand{\headrulewidth}{0.4pt}
\renewcommand{\footrulewidth}{0pt}

%-------------------------------------------------------------------------
% CHAPTER FORMATTING AND SPACING (via titlesec)
%--------------------------------------------------------------------------

\usepackage{titlesec}

% --------------------------------------------------------------------------
% CHAPTER - Nível 0 (Destaque Máximo)
% --------------------------------------------------------------------------
% Shape: display → standalone em nova linha
% Font: \Large\bfseries --- label e título ambos em bold para destaque máximo
% Numeração: SIM
% Espaçamento: 15pt antes, 28pt depois
% --------------------------------------------------------------------------

\titleformat{\chapter}[display]
{\normalfont\Large\bfseries}           % Formato: LARGE + bold (label)
{\chaptertitlename~\thechapter}        % Rótulo: "Chapter 1" (bold também)
{20pt}                                 % Espaço vertical antes do corpo
{\Large\bfseries}                      % Corpo do título: LARGE + bold

\titlespacing{\chapter}
{0pt}       % Margem esquerda (sem indent)
{15pt}      % Espaço ANTES do título do capítulo
{28pt}      % Espaço DEPOIS do título do capítulo (AUMENTADO)
[0pt]       % Margem direita

% --------------------------------------------------------------------------
% SECTION - Nível 1
% --------------------------------------------------------------------------
% Shape: hang → label à esquerda, corpo indentado (compacto)
% Font: \large\bfseries (mesmo destaque que chapter para coerência)
% Numeração: SIM
% Espaçamento: 15pt antes, 8pt depois (REDUZIDO de 15pt/10pt anteriores)
% --------------------------------------------------------------------------

\titleformat{\section}[hang]
{\normalfont\large\bfseries}           % Formato: large + bold
{\thesection}                          % Rótulo: "1", "2", etc
{1em}                                  % Espaço entre rótulo e corpo
{}                                     % Corpo do título (herda do format)

\titlespacing{\section}
{0pt}       % Margem esquerda (sem indent)
{15pt}      % Espaço ANTES do título da seção (REDUZIDO)
{8pt}       % Espaço DEPOIS do título da seção (REDUZIDO)
[0pt]       % Margem direita

% --------------------------------------------------------------------------
% SUBSECTION - Nível 2
% --------------------------------------------------------------------------
% Shape: hang → label esquerda, corpo indentado
% Font: \normalsize\bfseries (redução visual, ainda robusto)
% Numeração: SIM
% Espaçamento: 12pt antes, 6pt depois (compacto)
% --------------------------------------------------------------------------

\titleformat{\subsection}[hang]
{\normalfont\normalsize\bfseries}      % Formato: tamanho normal + bold
{\thesubsection}                       % Rótulo: "1.1", "1.2", etc
{1em}                                  % Espaço entre rótulo e corpo
{}                                     % Corpo do título

\titlespacing{\subsection}
{0pt}       % Margem esquerda (sem indent)
{12pt}      % Espaço ANTES do título da subseção
{6pt}       % Espaço DEPOIS do título da subseção
[0pt]       % Margem direita

% --------------------------------------------------------------------------
% SUBSUBSECTION - Nível 3
% --------------------------------------------------------------------------
% Shape: hang → label esquerda, máximo recuo (bem compacto)
% Font: \normalsize\bfseries (mesmo tamanho que subsection, padrão)
% Numeração: SIM (configurable com secnumdepth)
% Espaçamento: 10pt antes, 4pt depois (bem comprimido)
% --------------------------------------------------------------------------

\titleformat{\subsubsection}[hang]
{\normalfont\normalsize\bfseries}      % Formato: tamanho normal + bold
{\thesubsubsection}                    % Rótulo: "1.1.1", "1.1.2", etc
{1em}                                  % Espaço entre rótulo e corpo
{}                                     % Corpo do título

\titlespacing{\subsubsection}
{0pt}       % Margem esquerda (sem indent)
{10pt}       % Espaço ANTES do título
{4pt}       % Espaço DEPOIS do título (bem comprimido)
[0pt]       % Margem direita

% --------------------------------------------------------------------------
% PARAGRAPH - Nível 4 (CRÍTICO: Ilusão Óptica Resolvida)
% --------------------------------------------------------------------------
% Shape: runin → label e corpo na MESMA linha (inline, não block)
% Font: \normalsize\bfseries (tamanho normal + bold)
% Numeração: NÃO (secnumdepth=3 → paragraph não numerado)
% Espaçamento: 8pt ANTES do título (ADICIONADO) + 1em DEPOIS do título
% --------------------------------------------------------------------------

\titleformat{\paragraph}[runin]
{\normalfont\normalsize\bfseries}      % Formato: tamanho normal + bold
{}                                     % Rótulo: VAZIO (sem numeração)
{0em}                                  % Espaço entre rótulo e corpo (zero, runin)
{}                                     % Corpo do título

\titlespacing{\paragraph}
{0pt}       % Margem esquerda (sem indent)
{8pt}       % Espaço ANTES do título (ADICIONADO)
{1em}       % Espaço DEPOIS do título (ADICIONADO)
[0pt]       % Margem direita

% --------------------------------------------------------------------------
% SUBPARAGRAPH - Nível 5 (CRÍTICO: Ilusão Óptica Resolvida)
% --------------------------------------------------------------------------
% Shape: runin → label e corpo na MESMA linha (inline, não block)
% Font: \normalsize\itshape (tamanho normal + itálico, diferente de paragraph)
% Numeração: NÃO (secnumdepth=3 → subparagraph não numerado)
% Espaçamento: 8pt ANTES do título (ADICIONADO) + 1em DEPOIS do título
% --------------------------------------------------------------------------

\titleformat{\subparagraph}[runin]
{\normalfont\normalsize\itshape}       % Formato: tamanho normal + itálico
{}                                     % Rótulo: VAZIO (sem numeração)
{0em}                                  % Espaço entre rótulo e corpo (zero, runin)
{}                                     % Corpo do título

\titlespacing{\subparagraph}
{0pt}       % Margem esquerda (sem indent)
{8pt}       % Espaço ANTES do título (ADICIONADO)
{1em}       % Espaço DEPOIS do título (ADICIONADO)
[0pt]       % Margem direita

%--------------------------------------------------------------------------
% TABLES AND MACROS
%--------------------------------------------------------------------------

\usepackage{booktabs}
\usepackage{array}
\usepackage{longtable}
\usepackage{tabularx}

% Custom Columns
\newcolumntype{L}[1]{>{\raggedright\arraybackslash}p{#1}}
\newcolumntype{C}[1]{>{\centering\arraybackslash}p{#1}}
\newcolumntype{R}[1]{>{\raggedleft\arraybackslash}p{#1}}

% Macro for Visual Reference Links
\newcommand{\imglink}[1]{\hspace{2pt}\hyperref[#1]{\scriptsize\textbf{[Fig]}}}

%============================================================================
% HOTAS table environment (per Briefing v0.2.0.1)
%============================================================================

\newenvironment{hotastable}[1]{%
  \small
  \setlength{\tabcolsep}{2pt}
  \renewcommand{\arraystretch}{1.25}
  \begin{longtable}{L{1.00cm} L{0.90cm} L{0.90cm} L{3.30cm} L{6.40cm} L{1.40cm} L{2.10cm}}
  	\caption{#1}\\
  	\rowcolor{headerblue}
  	\multicolumn{1}{>{\centering\arraybackslash}p{1.00cm}}{\textbf{\color{white}Mode}} &
  	\multicolumn{1}{>{\centering\arraybackslash}p{0.90cm}}{\textbf{\color{white}Dir.}} &
  	\multicolumn{1}{>{\centering\arraybackslash}p{0.90cm}}{\textbf{\color{white}Act.}} &
  	\multicolumn{1}{>{\centering\arraybackslash}p{3.30cm}}{\textbf{\color{white}Function}} &
  	\multicolumn{1}{>{\centering\arraybackslash}p{6.40cm}}{\textbf{\color{white}Effect / Nuance}} &
  	\multicolumn{1}{>{\centering\arraybackslash}p{1.40cm}}{\textbf{\color{white}Dash34}} &
  	\multicolumn{1}{>{\centering\arraybackslash}p{2.10cm}}{\textbf{\color{white}Train.}} \\
  	\endfirsthead
  	\rowcolor{headerblue}
  	\multicolumn{1}{>{\centering\arraybackslash}p{1.00cm}}{\textbf{\color{white}Mode}} &
  	\multicolumn{1}{>{\centering\arraybackslash}p{0.90cm}}{\textbf{\color{white}Dir.}} &
  	\multicolumn{1}{>{\centering\arraybackslash}p{0.90cm}}{\textbf{\color{white}Act.}} &
  	\multicolumn{1}{>{\centering\arraybackslash}p{3.30cm}}{\textbf{\color{white}Function}} &
  	\multicolumn{1}{>{\centering\arraybackslash}p{6.40cm}}{\textbf{\color{white}Effect / Nuance}} &
  	\multicolumn{1}{>{\centering\arraybackslash}p{1.40cm}}{\textbf{\color{white}Dash34}} &
  	\multicolumn{1}{>{\centering\arraybackslash}p{2.10cm}}{\textbf{\color{white}Train.}} \\
  	\endhead
  	\multicolumn{7}{r}{\small\emph{Continued on next page}}\\
  	\endfoot
  	\endlastfoot
  }{%
  \end{longtable}
  }

%--------------------------------------------------------------------------
% SIMPLE REFERENCE MACROS FOR BMS DOCS
%--------------------------------------------------------------------------

\providecommand{\dashref}[1]{Dash-34~\S~#1}
\providecommand{\dashone}[1]{Dash-1~\S~#1}
\providecommand{\trnref}[1]{TRN~#1}
\providecommand{\trnman}{BMS Training Manual 4.38.1}
\providecommand{\bmsver}{Falcon BMS~4.38.1}
\providecommand{\dashrefs}[1]{\textit{TO 1F-16CMAM-34-1-1}, Dash-34, sections \texttt{#1}}

%--------------------------------------------------------------------------
% VERSION CONTROL MACROS
%--------------------------------------------------------------------------

\newcommand{\docversion}{v0.x.x.x}
\newcommand{\docbuild}{WIP Chapter 2}
\newcommand{\docstartdate}{06 February 2026}
\newcommand{\docenddate}{TBD}
\newcommand{\chapterscompletedof}{C2 WIP}
\newcommand{\tablesfilledpct}{N/A}
\newcommand{\fulldocversion}{\docversion+\docbuild}

%--------------------------------------------------------------------------
% GRAPHICS
%--------------------------------------------------------------------------

\usepackage{graphicx}
\graphicspath{{fig/}}
\usepackage{float}

%--------------------------------------------------------------------------
% TITLE
%--------------------------------------------------------------------------

\title{TMS, DMS and CMS Usage Guide for \bmsver}
\author{Carlos ``Metal'' Nader}
\date{Version \fulldocversion{} | Progress: Chapters \chapterscompletedof{} | Tables \tablesfilledpct{} | February 2026}

%============================================================================
% DOCUMENT BEGIN
%============================================================================

\begin{document}

\maketitle

\pagenumbering{roman}

%============================================================================
% TOC DEPTH CONFIGURATION (CRITICAL for report)
%============================================================================
\setcounter{tocdepth}{3}       % Show up to \subsubsection in TOC
\setcounter{secnumdepth}{3}    % Number up to \subsubsection

\newpage

\tableofcontents

\newpage

\pagenumbering{arabic}

%============================================================================
% WIP FILE METADATA (NOT RENDERED IN PDF)
%============================================================================

% File Name: chapter-C2-hotas-fundamentals-dev-2026-02-06.tex
% WIP Naming Convention: v1.4
% Target Chapter: C2 (HOTAS Fundamentals)
% Target Sections: 2.1, 2.2, 2.3
%
% WIP Status: dev
% Created: 2026-02-06
% Last Modified: 2026-02-06
% Integration Status: NOT YET INTEGRATED | INTEGRATION TARGET: v0.4.0.0 (MINOR bump)
%
% Narrative Completion: 100% (first draft)
% Table Fill Status: N/A (no tables in C2)
%
% Notes:
% - Chapter 2 provides conceptual foundation for C3 (DMS), C4 (TMS), C5 (CMS)
% - SOI content is conceptual only; operational tables remain in C3 (DMS)
% - Master modes coverage is minimal — only what's needed for HOTAS context
% - All three switches (DMS, TMS, CMS) presented in new pedagogical order
% - No tables, no diagrams — pure narrative
% - Estimated length: ~200-220 lines of content
% - Primary sources: Dash-34 § 2.1.1.2 (Moding and Sensor Concepts), § 2.1.5 (HOTAS), § 2.1.6 (MFDS)

%============================================================================
%============================================================================
% CHAPTER 2 — HOTAS FUNDAMENTALS
%============================================================================
%============================================================================

\chapter{HOTAS Fundamentals}
\label{chap:C2}

The F-16's HOTAS controls enable the pilot to execute critical flight, sensor, and weapon functions without removing hands from the throttle and stick grips. Within this architecture, three switches — the Display Management Switch (DMS), Target Management Switch (TMS), and Countermeasures Management Switch (CMS) — provide the primary interface for display management, sensor control, and defensive systems and form the core of the HOTAS architecture.

Effective use of these switches requires familiarity with two foundational concepts. The first is the Sensor of Interest (SOI), the mechanism that determines which display receives HOTAS inputs. The second is Master Modes (NAV, A-A, A-G), the operational contexts that govern available weapon systems and sensor configurations. This chapter establishes the conceptual framework necessary for the detailed switch 
behavior covered in Chapters~\ref{chap:C3}, \ref{chap:C4}, and \ref{chap:C5}.

%============================================================================
% SECTION 2.1 — SOI CONCEPT AND ARCHITECTURE
%============================================================================

\section{Sensor of Interest (SOI) — Concept and Architecture}
\label{sec:C2-S1}

The Sensor of Interest (SOI) is a designation that identifies which display currently receives HOTAS cursor slew commands and Target Management Switch (TMS) inputs. At any moment, only one display can hold SOI designation. This single-point control architecture enables the pilot to interact with multiple sensors and displays using HOTAS controls without manually selecting different switches for different sensors.

%--------------------------------------------------------------------------
% SUBSECTION 2.1.1 — THE SOI CONCEPT
%--------------------------------------------------------------------------

\subsection{The SOI Concept}
\label{sec:C2-S1-S1}

The F-16 cockpit contains three basic displays: the Head-Up Display (HUD) and two Multifunction Displays (MFD). The MFD can present various formats or pages, including radar, targeting pod, weapon management and tactical information. The pilot interacts with these displays through HOTAS controls, particularly the cursor slew control (CURSOR/ENABLE) and the Target Management Switch (TMS).

The Sensor of Interest (SOI) is the designation that identifies which display currently receives HOTAS inputs. At any moment, only one display holds SOI designation. When the pilot slews the cursor or presses TMS Up, the command is directed to whichever display is designated as SOI.

The SOI mechanism manages \textit{where} HOTAS inputs are directed, not \textit{what} those inputs do. The specific function of a HOTAS input — for example, whether TMS Up designates a target, breaks a track, or cycles through a mode — depends on the sensor mode, current Master Mode, and weapon selected. The SOI ensures that the input reaches the intended sensor; the sensor determines how the input is interpreted.

This design reflects the HOTAS philosophy: minimize pilot workload by consolidating control inputs through a limited number of switches while preserving the flexibility to interact with any sensor as needed.

%--------------------------------------------------------------------------
% SUBSECTION 2.1.2 — SOI INDICATORS AND VISUAL FEEDBACK
%--------------------------------------------------------------------------

\subsection{SOI Indicators and Visual Feedback}
\label{sec:C2-S1-S2}

The pilot can immediately identify which display holds SOI designation through unambiguous visual indicators:

\begin{itemize}
	\item \textbf{HUD as SOI:} An asterisk symbol (\texttt{*}) appears in the upper left corner of the HUD, above the airspeed scale.
	\item \textbf{MFD as SOI:} A border outline is drawn around the edges of the MFD display, forming a box that clearly distinguishes the SOI display from the non-SOI display.
	\item \textbf{NOT SOI indication:} When a display format is capable of being SOI but is not currently designated as such, the text \texttt{NOT SOI} may appear on the format. This text appears only on SOI-capable sensor formats, not on formats that are inherently unable to be SOI.
\end{itemize}

These visual cues allow the pilot to confirm at a glance which display is receiving HOTAS inputs. This immediate feedback is essential during dynamic operations, where the pilot may rapidly transition SOI between displays to manage multiple sensors or respond to tactical developments.

%--------------------------------------------------------------------------
% SUBSECTION 2.1.3 — SOI-CAPABLE DISPLAYS AND CONTEXT DEPENDENCY
%--------------------------------------------------------------------------

\subsection{SOI-Capable Displays and Context Dependency}
\label{sec:C2-S1-S3}

Not all displays can be designated as SOI, and the availability of SOI designation varies by master mode and display format. The displays capable of being SOI include the HUD, HMCS, and certain MFD formats: Fire Control Radar (FCR), Targeting Pod (TGP), Weapons (WPN), HARM Attack Display (HAD), and Horizontal Situation Display (HSD). Other MFD formats — such as Stores Management System (SMS), Data Transfer Equipment (DTE), TEST, and BLANK formats — are never SOI-capable, as they provide informational or configuration functions rather than sensor control.

The availability of SOI designation is further constrained by master mode. In Navigation (NAV) and Air-to-Ground (A-G) master modes, the HUD can be designated as SOI, allowing the pilot to manage HUD symbology and targeting cues directly through HOTAS inputs. In contrast, Air-to-Air (A-A), Dogfight (DGFT), and Missile Override (MSL OVRD) master modes restrict SOI designation to MFD formats only. The HUD cannot be designated as SOI in these air-to-air modes, ensuring that the pilot's primary sensor displays — typically the FCR or HSD — remain the focus of HOTAS control during air-to-air engagements.

This context-dependent availability reflects a deliberate design choice: master mode determines the operational context, and SOI availability is tailored to align the pilot's attention with the mission phase. Air-to-air modes prioritize radar and tactical displays; air-to-ground and navigation modes allow broader display flexibility, including HUD designation.

Understanding which displays can be SOI in which master modes is essential for effective HOTAS use. The pilot must know not only how to designate a display as SOI (covered in Chapter~\ref{chap:C3} on the Display Management Switch), but also whether such designation is permitted in the current operational context.

For complete operational reference tables showing SOI availability by master mode and the full list of valid and invalid SOI formats, see Chapter~\ref{chap:C3}, Section~\ref{sec:C3-S1}.

%============================================================================
% SECTION 2.2 — MASTER MODES AND CONTEXT-SENSITIVE BEHAVIOR
%============================================================================

\section{Master Modes and Context-Sensitive Behavior}
\label{sec:C2-S2}

Master modes are the highest-level configuration state of the F-16 avionics system. Each master mode pre-configures the cockpit displays, sensor modes, and weapon options for a particular mission phase. The master mode determines which HOTAS functions are available, which displays can be designated as SOI, and how individual switch inputs are interpreted by the avionics suite. Understanding master modes is essential to understanding why HOTAS switches behave differently in different operational contexts.

%--------------------------------------------------------------------------
% SUBSECTION 2.2.1 — MASTER MODE OVERVIEW
%--------------------------------------------------------------------------

\subsection{Master Mode Overview}
\label{sec:C2-S2-S1}

The F-16 provides four primary master modes, each optimized for a specific phase of flight or mission objective:

\begin{itemize}
	\item \textbf{Navigation (NAV):} The default master mode, used for transit, fuel management, and general situational awareness. In NAV, the avionics system prioritizes navigation displays and provides access to both air-to-air and air-to-ground sensors without committing to a specific weapon employment mode.
	
	\item \textbf{Air-to-Air (A-A):} Configures the avionics suite for beyond-visual-range (BVR) and within-visual-range (WVR) air-to-air engagements. A-A mode enables radar search and track modes, datalink integration, and air-to-air weapon selection. HUD symbology shifts to air-to-air cues, and the MFD formats default to radar and tactical displays.
	
	\item \textbf{Air-to-Ground (A-G):} Optimizes the cockpit for air-to-ground weapon delivery. A-G mode provides access to ground attack sensors (radar ground modes, targeting pod, HARM attack display), weapon delivery submodes (CCIP, CCRP, DTOS, etc.), and stores management functions. The HUD displays air-to-ground weapon delivery symbology.
	
	\item \textbf{Dogfight (DGFT) and Missile Override (MSL OVRD):} These are override modes accessed via a throttle-mounted switch. DGFT immediately configures the cockpit for close-range visual combat, typically defaulting to a high-off-boresight radar mode (ACM) and short-range missile (AIM-9) selection. MSL OVRD provides rapid access to medium-range missile employment (AIM-120) with optimized radar modes. Both override modes take precedence over any other master mode except Emergency Jettison.
\end{itemize}

Each master mode establishes default configurations for sensors, displays, and weapons. These defaults can be pre-programmed via the Data Transfer Cartridge (DTC) during mission planning or adjusted manually during flight. The master mode determines not only what the pilot sees on the displays, but also how HOTAS switches function. A single HOTAS input — such as TMS Up — produces different results in NAV, A-A, and A-G modes because the underlying sensor and weapon context has changed.

%--------------------------------------------------------------------------
% SUBSECTION 2.2.2 — MASTER MODE SELECTION
%--------------------------------------------------------------------------

\subsection{Master Mode Selection}
\label{sec:C2-S2-S2}

Master mode selection is designed for rapid reconfiguration with minimal pilot workload. The pilot can transition between master modes using dedicated switches on the Integrated Control Panel (ICP) or the throttle grip.

\textbf{Navigation (NAV)} is the default master mode and requires no explicit selection. If no other master mode is active, the avionics system operates in NAV. The pilot can return to NAV from A-A or A-G by ensuring the DGFT/MSL OVRD switch on the throttle is in the center (neutral) position and that no other master mode button is pressed.

\textbf{Air-to-Air (A-A)} is selected by pressing the \textbf{A-A button} on the ICP. Once selected, A-A mode remains active until the pilot selects a different master mode or returns the DGFT/MSL OVRD switch to center after an override mode engagement.

\textbf{Air-to-Ground (A-G)} is selected by pressing the \textbf{A-G button} on the ICP. As with A-A, the mode persists until changed by the pilot.

\textbf{Dogfight (DGFT)} is selected by moving the \textbf{DGFT/MSL OVRD switch} on the throttle grip to the \textit{left (outward)} position. This immediately overrides the current master mode and reconfigures the cockpit for visual combat. When the switch is returned to the center position, the avionics system reverts to the previously selected master mode (A-A, A-G, or NAV).

\textbf{Missile Override (MSL OVRD)} is selected by moving the \textbf{DGFT/MSL OVRD switch} to the \textit{right (inward)} position. This provides immediate access to medium-range missile employment. Returning the switch to center restores the previous master mode.

Table~\ref{tab:C2-master-mode-selection} summarizes the master mode selection controls.

\small
\renewcommand{\arraystretch}{1.2}

\begin{longtable}{L{4.5cm} L{4.0cm} L{5.5cm}}
	\caption{Master Mode Selection\label{tab:C2-master-mode-selection}}\\
	
	\rowcolor{headerblue}
	\textbf{\color{white}Master Mode} & \textbf{\color{white}Switch Location} & \textbf{\color{white}Switch Position / Action} \\
	\endfirsthead
	
	\rowcolor{headerblue}
	\textbf{\color{white}Master Mode} & \textbf{\color{white}Switch Location} & \textbf{\color{white}Switch Position / Action} \\
	\endhead
	
	\multicolumn{3}{r}{\emph{Continues on next page}}\\
	\endfoot
	
	\endlastfoot
	
	NAV (Navigation) & Default & No explicit selection required \\
	A-A (Air-to-Air) & ICP & Press A-A button \\
	A-G (Air-to-Ground) & ICP & Press A-G button \\
	DGFT (Dogfight) & Throttle & DGFT/MSL OVRD switch left (outward) \\
	MSL OVRD (Missile Override) & Throttle & DGFT/MSL OVRD switch right (inward) \\
\end{longtable}

\normalsize
\renewcommand{\arraystretch}{1.0}

The DGFT and MSL OVRD modes take precedence over A-A, A-G, and NAV, but do not override Emergency Jettison. This priority structure ensures that time-critical defensive or offensive actions can be executed immediately without navigating through multiple menu selections or button presses.

%--------------------------------------------------------------------------
% SUBSECTION 2.2.3 — WHY HOTAS SWITCHES BEHAVE DIFFERENTLY BY MASTER MODE
%--------------------------------------------------------------------------

\subsection{Why HOTAS Switches Behave Differently by Master Mode}
\label{sec:C2-S2-S3}

The context-sensitive design of HOTAS switches is a fundamental characteristic of the F-16 cockpit architecture. The same physical switch action produces different results depending on the master mode, the sensor mode, and the weapon selected. This design philosophy prioritizes pilot workload reduction: rather than requiring the pilot to select different switches for different tasks, the same switches adapt their behavior based on the current operational context.

Consider the Display Management Switch (DMS) Up command. In Navigation or Air-to-Ground master modes, pressing DMS Up designates the HUD as the Sensor of Interest, allowing the pilot to manage HUD symbology and targeting cues through HOTAS inputs. In Air-to-Air master modes, however, DMS Up has no effect — the HUD cannot be designated as SOI in A-A operations. The master mode determines whether the command is valid.

Similarly, the Target Management Switch (TMS) behavior varies by both master mode and which display is SOI. When the Fire Control Radar (FCR) is SOI in Air-to-Air mode, TMS Up commands a target designation, locking the radar onto a search contact. When the Targeting Pod (TGP) is SOI in Air-to-Ground mode, TMS Up may break an existing track or establish a cursor zero point, depending on the A-G submode. The physical switch action is identical, but the operational context — master mode, sensor mode, and SOI designation — determines the function executed.

This context-sensitivity extends to all three primary HOTAS switches covered in this guide. The DMS manages SOI selection and MFD format cycling, but the availability of certain SOI designations and format options depends on master mode. The TMS executes target management and designation functions, but those functions are interpreted differently by the FCR, TGP, WPN page, and other sensors. The Countermeasures Management Switch (CMS) operates relatively independently of SOI, but its interaction with the Countermeasures Dispenser System (CMDS) and Electronic Countermeasures (ECM) pod still varies by operational context and threat environment.

The pilot must be aware of the current master mode to predict HOTAS behavior. This awareness is not an additional cognitive burden — the master mode is a natural reflection of the mission phase. A pilot in Air-to-Air mode expects air-to-air HOTAS functions; a pilot in Air-to-Ground mode expects air-to-ground HOTAS functions. The master mode provides the contextual framework within which HOTAS inputs are interpreted, enabling a small number of switches to perform a wide range of mission-critical tasks.

Understanding this principle is essential for effective use of the detailed HOTAS tables and procedures in Chapters~\ref{chap:C3}, \ref{chap:C4}, and \ref{chap:C5}. Those chapters document the specific functions of DMS, TMS, and CMS in various master modes, sensor modes, and weapon configurations. The conceptual foundation established here — that master mode determines context, and context determines function — is the organizing principle behind those detailed references.

%============================================================================
% SECTION 2.3 — THE THREE SWITCHES: DMS, TMS, AND CMS
%============================================================================

\section{The Three Switches: DMS, TMS, and CMS}
\label{sec:C2-S3}

The Display Management Switch (DMS), Target Management Switch (TMS), and Countermeasures Management Switch (CMS) form the core of HOTAS sensor and display management in the F-16. Each switch serves a distinct function, and understanding their roles in relation to SOI and master modes is essential before examining the detailed behaviors covered in subsequent chapters.

%--------------------------------------------------------------------------
% SUBSECTION 2.3.1 — DISPLAY MANAGEMENT SWITCH (DMS)
%--------------------------------------------------------------------------

\subsection{Display Management Switch (DMS) — Overview}
\label{sec:C2-S3-S1}

The Display Management Switch (DMS) is a four-way, spring-loaded hat switch located on the flight stick. Its primary functions are SOI selection and MFD format cycling. The DMS does not perform tactical actions such as target designation or weapon release; instead, it manages \textit{which display or sensor} receives HOTAS inputs from other controls.

DMS Up attempts to designate the HUD as the Sensor of Interest, provided the current master mode permits HUD SOI designation (NAV and A-G allow this; A-A, DGFT, and MSL OVRD do not). DMS Down cycles the SOI designation between the left and right MFDs, or transitions SOI from the HUD back to an MFD. DMS Left and DMS Right cycle through the MFD format pages displayed on the left and right MFDs, respectively, independently of which display is currently SOI.

The DMS is a transversal SOI manager. It controls the routing of HOTAS inputs to displays, but does not execute sensor-specific or weapon-specific functions. This distinction separates the DMS from the TMS, which operates \textit{through} the SOI to execute tactical actions such as target designation, track management, and weapon cueing.

Detailed DMS functionality, including SOI selection mechanics, MFD format cycling behavior, and master mode constraints, is covered in Chapter~\ref{chap:C3}.

%--------------------------------------------------------------------------
% SUBSECTION 2.3.2 — TARGET MANAGEMENT SWITCH (TMS)
%--------------------------------------------------------------------------

\subsection{Target Management Switch (TMS) — Overview}
\label{sec:C2-S3-S2}

The Target Management Switch (TMS) is a four-way, spring-loaded hat switch located on the throttle grip. Its primary function is target management and designation, controlling how the pilot interacts with radar contacts, targeting pod tracks, weapon employment cues, and other sensor-derived targeting information.

Unlike the DMS, which manages \textit{which} display is active, the TMS operates \textit{through} the currently active SOI to execute tactical functions. The specific behavior of each TMS direction depends on the master mode, the sensor mode, and which display holds the SOI designation. For example, when the Fire Control Radar (FCR) is SOI in Air-to-Air mode, TMS Up commands a target designation, locking the radar onto a search contact. When the Targeting Pod (TGP) is SOI in Air-to-Ground mode, TMS Up may break an existing track, establish a cursor zero point, or transition to a post-designate mode, depending on the A-G submode and weapon selected.

The TMS is fundamentally context-dependent. The pilot must understand not only the current master mode, but also which display is SOI and what sensor mode that display is operating in. This context-sensitivity allows a single four-way switch to perform dozens of distinct functions across the range of F-16 mission profiles.

Detailed TMS functionality, including target designation procedures, sensor-specific behaviors, and master mode contexts, is covered in Chapter~\ref{chap:C4}.

%--------------------------------------------------------------------------
% SUBSECTION 2.3.3 — COUNTERMEASURES MANAGEMENT SWITCH (CMS)
%--------------------------------------------------------------------------

\subsection{Countermeasures Management Switch (CMS) — Overview}
\label{sec:C2-S3-S3}

The Countermeasures Management Switch (CMS) is a four-way, spring-loaded hat switch located on the throttle grip. Its primary function is countermeasures dispensing, controlling the release of chaff and flare from the ALE-47 Countermeasures Dispenser System (CMDS) and managing Electronic Countermeasures (ECM) pod transmissions.

Unlike the DMS and TMS, the CMS operates largely independently of Sensor of Interest designation. CMS commands interact directly with defensive systems — the CMDS, ECM pod, and Radar Warning Receiver (RWR) — rather than routing through the SOI architecture. The pilot actuates the CMS in response to threat warnings displayed on the RWR, and the countermeasures system responds according to pre-programmed or manually configured dispense patterns.

CMS Up, CMS Down, CMS Forward, and CMS Aft each trigger specific countermeasures responses. The exact chaff and flare quantities dispensed, and the enabling or disabling of ECM transmissions, depend on the CMDS programming and the CMS manual mode settings (if selected). The CMS behavior is relatively consistent across master modes, although certain operational contexts (such as air-to-ground attack profiles) may favor specific countermeasures programs.

The CMS provides the pilot with immediate, hands-on control of defensive systems without requiring menu navigation or display interaction. This direct-access design reflects the time-critical nature of countermeasures employment in response to surface-to-air or air-to-air threats.

Detailed CMS functionality, including CMDS program interaction, ECM pod control, and manual mode operation, is covered in Chapter~\ref{chap:C5}.

%============================================================================
% END OF CHAPTER 2
%============================================================================

\end{document}
